\chapter{Conclusion}
\label{ch:conclusion}
\newpage

\section{Summary of the Work}
\label{sec:ch5-summary}

This thesis asked whether the code that autonomous coding agents submit to real projects follows the
SOLID design principles. Chapter~\ref{ch:introduction} motivated the question with the distance
between benchmark scores and acceptance by maintainers. Chapter~\ref{ch:literature} showed that
existing studies either measure design quality on small sets of tasks or do not compare agents on a
common corpus of real contributions. Chapter~\ref{ch:methodology} defined one deterministic detector
per principle, validated against an independent reading judge. Chapter~\ref{ch:results} applied the
design to 4{,}245 analysable Java pull requests from 1{,}010 repositories, written by three agents
and including merged, rejected and open requests, and repeated it on 4{,}946 analysable Python pull
requests from 2{,}921 repositories.

Three design decisions distinguish the study. The corpus is not limited to merged code, so the effect
of review can be examined. The validity of each detector is measured: 1{,}495 Java and 980 Python
entities were read and labelled with a reason for each label. And every Java result was tested again
in a language without Java's structural guarantees, and only results found in both languages are
called general.

This chapter answers the research questions (Section~\ref{sec:rq-answers}), lists the contributions
(Section~\ref{sec:contributions}) and practical implications (Section~\ref{sec:practical}), states
the limitations (Section~\ref{sec:limitations}) and proposes future work
(Section~\ref{sec:future-work}).

\section{Answers to the Research Questions}
\label{sec:rq-answers}

RQ1 to RQ3 are answered for Java, where they were first established. RQ4 states which answers the
Python replication supports.

\paragraph{The short answer.} Agents do violate SOLID, and they do so rarely. After correcting for
measured precision, between 0.4\% and 4.1\% of Java pull requests and between 0.3\% and 3.5\% of
Python pull requests introduce a violation of a given principle. The same code already contains
7{,}501 flagged classes, 2{,}343 flagged dispatch chains, 816 flagged contracts and 183
refused-bequest relationships in Java, against 16 pull requests that introduced a refused bequest.
The rate at which agents create structural problems is therefore low, and the quantity of structural
problems in the code they work in is high. Which agent made the change matters less than how large
the change was: each doubling of the files touched multiplies the odds of a flag by 1.88 to 1.95 in
Java and 2.01 to 2.18 in Python, and only one agent difference survives that adjustment in both
languages.

The detectors agree with the reading judge for two principles and screen for three. Dependency
Inversion reaches precision 0.892 in Java on a census of all 195 flagged pull requests and 1.000 in
Python, and the Liskov flag restricted to refused bequest reaches 1.000 on 64 Java flags across a
fitted sample and a holdout, and 0.857 in Python. Single Responsibility, Open/Closed and Interface
Segregation have precision between 0.110 and 0.425, so their flags locate code worth reading and
should not be counted as violations. Against published analysers the Single Responsibility rubric
agrees with PMD at $\kappa = 0.574$ and with DesigniteJava at $\kappa = 0.531$, while those two tools
agree with each other at $\kappa = 0.384$; the Open/Closed detector is the one that no published tool
corroborates, at $\kappa = 0.122$.

\paragraph{RQ1: do agents differ, where, and does review select on it?}
The agents differ less than pooled rates suggest. Claude Code has the highest pooled flag rate on all
five principles and OpenAI Codex the lowest on four; Codex introduced no refused-bequest violation
in 770 pull requests (Table~\ref{tab:agent-rates}).

Much of this ordering reflects change size. Claude Code touches 46.4 Java files per pull request on
average and Codex 10.9, and flag rates rise with size for every agent, which matches the churn effect
reported by \citet{cynthia2026postmerge}. Within size bands (Table~\ref{tab:size-control}) the
ordering holds for the largest changes, where Claude Code flags 31.3\% of its Open/Closed changes
against 8.9\% for Codex, but not for the smallest, where Copilot is highest on two principles and
Codex on Dependency Inversion.

The joint model (Table~\ref{tab:mixed-model}) accounts for size, time and repository together. Each
doubling of the files touched multiplies the odds of a flag by about 1.9 on every modelled principle.
One agent difference remains: Claude Code's Open/Closed odds are 2.22 times Copilot's (1.21--4.07).
The Interface Segregation and Dependency Inversion differences fall to 1.11 and 1.29, with intervals
that include 1. On its own, the Java Open/Closed coefficient does not pass a Bonferroni correction
over the six agent coefficients; it is accepted because it replicates in Python. Time has no
detectable effect (all $p > 0.5$). The pooled and macro rates disagree for Codex on two principles,
and Codex has a third as many pull requests as Copilot, so its estimates are the least precise.

On outcomes, with the reconstruction method held constant, merged pull requests are structurally
indistinguishable from those closed without review, which are 1{,}006 of the 1{,}177 unmerged
requests (Table~\ref{tab:outcome-stratum}). The two strata that record a maintainer's judgement are
small (36 and 135 requests) and have higher Interface Segregation and Dependency Inversion rates than
merged code, so whether review acts on structure remains open. For these signals, a corpus of merged
code is not a biased sample of unreviewed agent code.

\paragraph{RQ2: can SOLID be measured at scale, and how faithfully?}
Partly, and the answer differs by principle. Dependency Inversion reaches precision 0.892, measured
on every flagged pull request, while the other four detectors have precision between 0.330 and
0.420. Weighted recall is 0.742 to 1.000 for SRP, OCP and DIP, whose flag rates can therefore be read
as upper bounds, but 0.386 for Liskov and 0.464 for Interface Segregation
(Table~\ref{tab:validation}).

Two detectors improve with a restriction suggested by the labels. Restricting the Liskov flag to
refused bequest gives precision 1.000 on 64 flags over two samples. Requiring interface contracts
of at least 30 members raises Interface Segregation precision to 0.786 pooled. On a holdout of 200
new entities per principle, the Liskov restriction held (33 of 33) and the width condition fell from
0.862 to 0.704, outside its own holdout interval. The categorical rule kept its precision; the
numeric threshold did not.

The remaining gap reflects the difference between structure and intent. Open/Closed precision is flat
across dispatch chain lengths, because whether a set of variants can grow does not depend on how many
branches it currently has.

The detectors were also compared with three published analysers on 111{,}029 production types. The
Single Responsibility rubric agrees with PMD at $\kappa = 0.574$ and with DesigniteJava at
$\kappa = 0.531$, while PMD and DesigniteJava agree at $\kappa = 0.384$ (Table~\ref{tab:tri-tool}),
and 68.0\% of the rubric's flags are also flagged by at least one of the tools. Where the tools
diverge, the reason is the unit of measurement: DesigniteJava assigns a substitutability fault to the
supertype and cycles to packages, while this study assigns them to the subclass and to compilation
units.

\paragraph{RQ3: where does risk concentrate?}
Most flagged problems in the touched code were there before the agent's change. After precision
correction, the Java code these agents touched contains about 2,600 Single Responsibility violations
and 183 refused-bequest relationships, and between 0.4\% and 4.1\% of pull requests introduced a
violation of a given principle (Table~\ref{tab:present-introduced}). For refused bequest, 16 pull
requests introduced one.

The reason is structural. Single Responsibility and Open/Closed violations lie inside one unit, which
a pull request can write in full. Liskov, Interface Segregation and Dependency Inversion violations
lie in relations between units, and one contribution rarely writes both ends. Review of single
changes can therefore catch problems inside units but will rarely see problems between units.

\paragraph{RQ4: which answers hold in Python?}
Most of them. With the same agents, curation rules, validation procedure and model, and thresholds
recalibrated on Python before any Python label existed, the Python study reproduces the main Java
results (Table~\ref{tab:cross-language}).

Change size dominates in both languages, with odds ratios of 2.01 to 2.18 per doubling in Python
against 1.88 to 1.95 in Java. The Open/Closed agent difference replicates: Claude Code's adjusted
odds are 2.09 times Copilot's in Python (1.44--3.04, passing the Bonferroni threshold), against 2.22
in Java, and its Python Open/Closed flags are the most often confirmed of the three agents. A
possible second difference, more import cycles introduced by Claude Code (1.88, 1.17--3.03), does not
pass the correction and is not supported by Java. Pre-existing problems again outnumber new ones: 176
refused-bequest relationships are present against five introducing pull requests, and between 0.3\%
and 3.5\% of Python pull requests introduce a genuine violation. Applied unchanged to the Python
labels, the categorical Liskov rule scores six of seven and the numeric width rule one of five. And
46 of the 51 confirmed Python Liskov violations rest on faults that \texttt{javac} would reject.

Two results do not transfer. The Single Responsibility screen degrades, from precision 0.350 to
0.110, because the large classes in Python projects are mostly shaped by frameworks. And the outcome
result differs: after adjusting for size and agent, merged Python code carries fewer Open/Closed and
Liskov flags than unmerged code (odds ratios 1.57 and 2.43 for unmerged requests), while merged and
unreviewed Java code were indistinguishable. Since most unmerged requests in both languages were never
reviewed, the Python result describes abandoned code, but it means that the Java finding about
merged-only corpora does not generalise.

\section{This Study Against the State of the Art}
\label{sec:against-sota}

Chapter~\ref{ch:literature} set out what earlier work reports. Table~\ref{tab:sota-comparison}
places each of those findings beside the corresponding measurement made here. Three of them are
confirmed on real pull requests, two are not tested by this design, and two point in a different
direction once the unit of measurement is made explicit.

\begin{table}[p]
\centering
\caption{Earlier findings and what this study measures against them. ``Not tested'' means the design
does not observe the quantity in question, not that the earlier finding is doubted.}
\label{tab:sota-comparison}
\footnotesize
\begin{tabular}{@{}>{\raggedright\arraybackslash}p{3.5cm}
                  >{\raggedright\arraybackslash}p{4.6cm}
                  >{\raggedright\arraybackslash}p{4.9cm}@{}}
\toprule
\textbf{Earlier finding} & \textbf{This study} & \textbf{Outcome} \\
\midrule
Code volume predicts structural degradation, the Volume-Quality Inverse Law
\citep{zhu2026smells} &
Change size is the dominant term in the joint model, OR 1.88--1.95 per doubling in Java and
2.01--2.18 in Python, on every modelled principle &
Confirmed, and it is the strongest result in this thesis \\
\addlinespace
Agent differences disappear once code churn is normalised \citep{cynthia2026postmerge} &
After adjusting for files touched, repository and time, only Claude Code on Open/Closed survives in
both languages (2.22 Java, 2.09 Python) &
Confirmed, with one exception that replicates \\
\addlinespace
Passing tests does not imply a change a maintainer would accept \citep{metr2026sweverified} &
Merged Java pull requests are structurally indistinguishable from those closed without review, and
only 171 Java and 196 Python unmerged requests were ever reviewed &
Consistent: merge status carries little structural information here \\
\addlinespace
Generated projects reach 91\% functional correctness yet carry thousands of design issues
\citep{beyond2026designissues} &
Violations already present in the touched code outnumber those introduced by a factor of about
eleven in Java and thirty-five in Python for refused bequest &
Consistent, but the unit differs: that study builds new projects, this one measures contributions to
existing ones \\
\addlinespace
Agents succeed on 4.3\% of dependency-control and 15.2\% of responsibility-decomposition tasks
\citep{zhu2026nitr} &
Agents seldom introduce dependency or responsibility violations into existing repositories &
Different question: that benchmark asks agents to repair structure, this study asks whether they
damage it \\
\addlinespace
Structural erosion grows across 77\% of long-horizon trajectories \citep{slopcodebench2026} &
The time coefficient is inert for every delta principle in both languages &
Not tested: the unit here is one pull request, not a trajectory of successive edits \\
\addlinespace
Single Responsibility is the easiest principle for a language model to detect and Dependency
Inversion the hardest \citep{pehlivan2025solid} &
For structural detectors the order is reversed: Dependency Inversion is the most precise
(0.892 Java, 1.000 Python), Single Responsibility the least (0.350 Java, 0.110 Python) &
Different construct: that work measures a model reading examples, this one measures a static rule
over a corpus \\
\bottomrule
\end{tabular}
\end{table}

Two of these entries deserve comment. The comparison with \citet{zhu2026nitr} is not a
disagreement. That benchmark gives an agent a maintainability task and records whether the agent
completes it, so a 4.3\% success rate says agents are poor at repairing dependency structure. This
study records whether an agent's change creates a violation, and finds that it seldom does. An agent
can be both unlikely to fix a cycle and unlikely to create one, and the two results together suggest
that agents leave dependency structure roughly as they found it.

The reversal against \citet{pehlivan2025solid} follows from what each method can see. A language
model reading a short example can recognise a class doing several things, which is why Single
Responsibility is easy for it, but it cannot see an import cycle that spans files it was not shown.
A static analyser over a whole repository has the opposite profile: the import graph is exactly what
it computes, while responsibility is a judgement about intent that no metric captures, which is why
the Single Responsibility rubric reaches precision 0.350 in Java and 0.110 in Python. The two
findings are compatible and, taken together, argue for combining a graph-based detector with a
reading model, not for choosing between them.

\section{Contributions}
\label{sec:contributions}

\begin{enumerate}[label=\textbf{C\arabic*.}, leftmargin=1.6cm]
  \item A repository-scale measurement instrument for the five SOLID principles, with each
  detector's precision and recall measured against an independent reading judge. None of the studies
  reviewed in Chapter~\ref{ch:literature} measures SOLID conformance on code that agents authored for
  real projects.
  \item A corpus of agent pull requests covering all three outcomes, which makes it possible to
  compare merged code with rejected and open code. Merged and unreviewed code are indistinguishable
  in Java and differ in Python.
  \item A description of which SOLID principles a structural detector measures reliably and which it
  can only screen, including two restrictions derived from labelled false positives, tested on a
  holdout and in a second language, and an explanation of why no width threshold can make the
  Open/Closed detector precise.
  \item Evidence, in two languages, that most design problems in the code agents touch were present
  before their changes, and that relation-level violations are rarely introduced by a single pull
  request.
  \item An observation about type systems and measurement: signals that the Java compiler rules out
  account for 46 of the 51 Liskov violations confirmed in Python, so the set of detectable violations
  depends on the language.
  \item An auditable validation procedure: every label must trace back to the source the judge was
  shown, every label has a recorded reason, and a mechanical audit enforces both.
  \item A replication with the same agents in a second language, which separates results that hold
  for agent-written code in both languages (the effect of change size, the Open/Closed difference,
  pre-existing violations, the behaviour of fitted rules) from results that do not (the outcome
  comparison, the precision of the size-based responsibility screen).
\end{enumerate}

\section{Practical Implications}
\label{sec:practical}

For a team choosing a coding agent, the SOLID profile is an addition to functional benchmarks, not a
replacement. The differences between agents are small and mostly reflect how much each agent changes
at once. The one difference found in both languages is Claude Code's higher rate of Open/Closed flags.
Across all three agents and both languages, classes the agents authored are flagged less often than
the existing classes they edit, at about half the rate in Java and three quarters in Python.

For a team building an automated review check, the results suggest one addition and one caution. The
check worth adding is refused bequest: a subclass that overrides a working inherited method with a
body that only refuses. It is cheap to compute, had precision 1.000 on 64 Java flags across two
samples and six of seven in Python, and describes something a reviewer would object to. In a
language without a compiler, a second check is that an override still accepts the calls its base
accepts, since such faults made up most of the confirmed Python Liskov violations. An interface-width
threshold does not transfer: it lost 16 points on the Java holdout and scored one of five in Python.
A Single Responsibility or Open/Closed flag from detectors of this kind should prompt a discussion,
not block a merge, because most such flags did not survive reading in either language.

Because the violations that matter most build up across changes, an automated structural check adds
most where review of single changes cannot look: across the history of a component.

\section{Limitations}
\label{sec:limitations}

\paragraph{SOLID adherence is not design quality.} The thesis measures conformity to five
principles; it does not measure whether conforming code is cheaper to change, less fault-prone or
easier to understand. No step observes such an outcome. The link from a violation to later cost comes
from the literature, which is not unanimous: \citet{olbrich2010godclass} find that God Classes are
less change-prone and fault-prone than other classes once size is controlled. A reader who does not
accept SOLID as a quality model can read the results as measurements of stylistic conformity; the
comparisons between agents remain valid under that reading because the same instrument is applied to
every agent. Linking these flags to later defects or changes in the same repositories would answer
the question.

\paragraph{Every agreement figure is agreement with an LLM judge.} Precision and recall measure
agreement between a deterministic detector and a language model that reads the source. LLM judgement
of SOLID violations has been evaluated by \citet{pehlivan2025solid}, but the judge can still be wrong,
including in the same direction as a detector. Two checks constrain this. For Java, 174 entities
judged twice agreed on 89.7\%; no Python entity was judged twice, so the Python judge's consistency
is not measured. And for Java the class-level detectors agree with three independent analysers at
least as closely as those analysers agree with each other. Interface Segregation has no counterpart in
those tools, and the Open/Closed detector agrees with DesigniteJava's type-conditional smells at only
$\kappa = 0.122$, so for these two principles, and for all Python figures, the judge is in practice
the only check.

\paragraph{The judge and one agent come from the same model family.} The judge is a Claude model and
Claude Code is one of the agents. Per-agent precision gives no sign of leniency towards Claude Code
(Table~\ref{tab:validation-agent}), but it cannot detect a bias shared by the judge and the detector.

\paragraph{Several agent coefficients are tested.} The joint models test six agent coefficients in
Java and ten in Python. With a Bonferroni correction, the Java Open/Closed coefficient
($p = 0.010$) would not pass alone and the Python import-cycle coefficient ($p = 0.0096$) does not
pass. Only the Open/Closed difference, which passes in Python and has the same direction and a
similar size in Java, is presented as a finding.

\paragraph{Change size is adjusted for, not explained.} The agents write changes of very different
sizes, and flag rates rise with size. The model adjusts for the number of files touched. If writing
larger changes is part of how an agent works, the adjustment also removes that part of its effect.
The model does not measure the kind of task each agent was given, which remains part of its
coefficient.

\paragraph{Time is modelled as a linear trend.} Median pull request dates differ by about seven months
between agents in Java and about two months in Python. The time coefficient is close to 1 for every
delta principle in both languages. A linear term cannot detect an abrupt change, such as a model
release that altered one agent's behaviour, so the agent coefficients are adjusted for gradual drift
only. The Python Single Responsibility trend rests on a flag with precision 0.110.

\begin{sloppypar}
\paragraph{Rejected does not mean refused on quality.} Most closed pull requests were never reviewed.
Only the two reviewed strata, \texttt{changes\_requested} and \texttt{reviewed\_then\_closed}, record a
judgement, and they are small: 36 and 135 requests in Java, 33 and 163 in Python. The conclusions
drawn from the outcome comparison are correspondingly narrow. In Java a merged-only corpus is not
biased relative to unreviewed code; in Python merged code is cleaner than abandoned code. Neither
shows what reviewers select on.
\end{sloppypar}

\paragraph{Reconstruction differs by outcome.} In Java, 62.1\% of merged requests are analysable
against 56.8\% of rejected ones, and the methods available to recover the before tree differ. The
outcome comparison is therefore made only within matched methods.

\paragraph{Symbol resolution is incomplete.} Median supertype resolution in Java is 87.0\%, and 718
pull requests resolve below 50\%, mostly in framework-heavy projects, where Liskov and Interface
Segregation undercount. An unresolved supertype removes candidates and does not create false ones,
so these rates are lower bounds. The Python extractor has no symbol solver and resolves a base class
only when it is defined in the repository, so classes that derive from framework types contribute no
override pairs or contracts; the error has the same direction.

\paragraph{One snapshot per pull request.} The design compares two trees. Gradual erosion over many
changes, which the locus-of-risk results suggest matters most for relation-level principles, is not
observed.

\paragraph{Two languages from one dataset.} Both corpora come from AIDev, so any property of how the
dataset selects agent pull requests affects both. Two languages also do not establish generality: both
are widely used, have large open-source ecosystems and use nominal class hierarchies, and a statically
typed language with structural typing is not represented. A benchmark built on specific pull requests
also decays: 18 Java requests could no longer be retrieved because repositories were deleted or their
history rewritten.

\paragraph{Python loses source unevenly.} 496 Python pull requests (7.0\%) could not be analysed
because their branch was deleted and neither a merge commit nor recorded commits remained. The loss is
11.0\% for rejected requests against 5.1\% for merged ones, and 14.3\% for Copilot against 3.6\% and
3.1\% for the other agents. The lost requests were not replaced, so the Python outcome and Copilot
figures describe the retrievable part of those populations.

\paragraph{The Python instrument is weaker in two places.} The Single Responsibility screen is correct
on 11 of 100 Python flags, so Python SRP rates describe large classes more than classes with several
responsibilities. The Dependency Inversion validation confirms that each flagged pull request closes
a new import cycle through changed code, but not that the cycle is harmful, and the most common form,
a module importing its own package root, is accepted in many projects.

\paragraph{The two languages define two principles differently.} Python has no \texttt{interface}, so
its Interface Segregation flag requires width together with low implementer coverage or a forced stub,
while the Java flag can fire on width alone; and each language's thresholds were calibrated on its own
corpus. Flag rates are therefore compared across languages only through the models and contrasts of
Section~\ref{sec:python}, never as raw rates. The fall from 6.5\% of Java pull requests to 0.6\% of
Python ones on Interface Segregation is largely a difference of definition.

\paragraph{A few Python classes are vendored.} 489 of 124{,}143 Python class rows lie inside a virtual
environment committed to a repository, 469 of them from one pull request, and 41 of the flagged SRP
classes are among them. The production filter does not exclude committed environments; the effect on
any pull-request-level rate is one pull request.

\paragraph{Small counts for relation-level principles.} Sixteen Java pull requests introduce a Liskov
violation and 195 a dependency cycle; in Python the counts are 47 for Liskov, 32 for Interface
Segregation and 173 for dependency cycles. The Liskov and Interface Segregation agent coefficients
have intervals several times wider than the Open/Closed one. Counts are reported next to every rate so
that readers can judge them directly.

\section{Future Work}
\label{sec:future-work}

\begin{itemize}
  \item \textbf{A change-history signal for Open/Closed.} Dispatch chain length does not predict
  Open/Closed violations. A conditional that is edited in many pull requests to add cases is, by
  definition, not closed for modification, and a history of changes could measure this directly.
  \item \textbf{A closed-set suppressor for external specifications.} The current suppressor
  recognises only enumerations and removed none of the flagged validation rows. The closed sets that
  occur are JDBC type codes, protobuf scalars, Unicode blocks and generated parser tables in Java, and
  vendor API event types, protocol method names and grammar node kinds in Python.
  \item \textbf{An Open/Closed proxy that agrees with published smells.} The comparison with
  DesigniteJava reported in Section~\ref{sec:crosscheck-results} reaches only $\kappa = 0.122$, and the
  detector fires three times as often. Requiring that a hierarchy already exist over the dispatched
  types, or that the chain sit in a class with no abstract supertype, would move the proxy towards the
  published definition and can be tested against the smells already recorded.
  \item \textbf{Recall for Liskov and Interface Segregation.} The weighted recall of these detectors
  depends on the few violations found in a large unflagged stratum. A sample drawn from unflagged
  entities only would narrow the intervals; the three Liskov misses found so far suggest looking at
  silent refusals.
  \item \textbf{Judge consistency in Python and a second judge.} Re-judging a Python sample, and
  judging a sample with a model from a different family, would measure the judge's stability and its
  independence from the evaluated agents.
  \item \textbf{Why merged and rejected code differ in Python but not in Java.} A corpus with more
  reviewed rejections in both languages, and an analysis of what reviewers wrote when closing a
  request, would show whether the Python difference comes from review or from abandonment.
  \item \textbf{A framework-aware responsibility screen for Python.} Most Python SRP false positives
  are classes whose size a framework dictates. Recognising a class by the framework base it derives
  from could recover the precision lost between the languages.
  \item \textbf{A production filter for implementers.} Most Python Interface Segregation false
  positives are test doubles, null objects and unfinished scaffolds that stub a contract on purpose.
  Excluding implementers under test paths would remove the first group automatically.
  \item \textbf{A third language with a different type system.} A statically typed language with
  structural typing would test the compiler observation again, since such a language rejects some
  substitutability faults and not others.
  \item \textbf{Repository history as the unit of analysis.} Relation-level problems build up outside
  the reach of any single change. Measuring a component across its history is the natural next design.
\end{itemize}

\section{Closing Remark}
\label{sec:ch5-closing}

The central result of this thesis is a distinction. Most design problems that the detectors find in
the code these agents touched were there before the agents arrived, and new violations appear mainly
inside single units, where one change can create them. Relation-level problems accumulate across many
changes and many authors. This held for three agents, for merged and unmerged work, in Java and again
in Python, as did the finding that how much an agent changes matters more than which agent it is. The
remaining question is what happens at a larger scale: a repository maintained over months by agents
and by the people who edit their output.
